\section{Introduction}\label{sec:introduction}

Matrix multiplication is a basic example of a bilinear map whose
algebraic structure permits fewer products than its entrywise formula
suggests. Strassen's seven-product algorithm for $2\times2$ matrices
established this possibility~\cite{strassen}. For $3\times3$ matrices,
the usual formula uses $27$ products, whereas Laderman's algorithm uses
$23$~\cite{laderman}. The minimum remains unknown. Its determination
asks both for efficient constructions and for structural obstructions
that apply to every possible bilinear algorithm.

A construction specifies linear forms in both inputs and the
coefficients combining their products into the output. Quotient-rank
bounds constrain a smaller set of data: the positions and
multiplicities of the first-input forms. In this paper we combine
these constraints with a split flattening of the complete
multiplication tensor. At a hypothetical length of twenty, the two
sets of conditions force equality in matrix-rank subadditivity.
The equality case then yields relations among all three factor lists,
strong enough to exclude the decomposition.

Write $\T$ for the $3\times3$ multiplication tensor and
$R_{\F}(\T)$ for its exact tensor rank over the two-element field $\F$.
This rank is the least number of products of a linear form in the first
input with a linear form in the second, with linear operations otherwise
free~\cite{bcs}. We prove the following statement without imposing symmetry or
support conditions on the algorithm.

\begin{theorem}\label{thm:main}
Every exact bilinear algorithm for $3\times3$ matrix multiplication over
$\F$ uses at least $21$ scalar multiplications. Equivalently,
\begin{equation}\label{eq:main-bound}
 R_{\F}(\T)\geq21.
\end{equation}
\end{theorem}

Together with the explicit $23$-term decomposition in
Appendix~\ref{app:upper}, this gives
$21\leq R_{\F}(\T)\leq23$.
Wang~\cite{wang21} obtained the same numerical lower bound by
capacity-constrained profile enumeration. The proof here identifies a
different obstruction: finite geometry forces a rank equality, and
that equality supplies cross-factor relations which contradict the
forced first-factor configuration.

\subsection{The equality mechanism}
Consider a decomposition
$\T=\sum_{t=1}^{r}A_t\otimes B_t\otimes C_t$, with the factors
represented by $3\times3$ coefficient matrices. For a subspace
$W\leq M_3(\F)$, quotienting the first tensor factor by $W$ removes
every term with $A_t\in W$. A lower bound for the quotient rank
therefore gives an upper bound on the number of such terms. We call
these \emph{occupation bounds}.

A second constraint retains the row and column indices of $A_t$.
Splitting these indices gives a $27\times27$ matrix $P$ for the
multiplication tensor and matrices $M_t$ for its summands. The matrix
$P$ is a permutation matrix, while $\rank M_t=\rank A_t$ for
nonzero factors. Hence
\begin{equation}\label{eq:intro-budget}
 P=\sum_tM_t,\qquad 27\leq\sum_t\rank A_t.
\end{equation}
This lower bound on the sum of matrix ranks is not by itself a lower
bound of $21$ on the number of terms. Its role is to combine with
occupation geometry.

At length $20$, eight strengthened quotient bounds force the
first factors of rank at least two to have pairwise rank-one
differences. Such matrices lie in a common affine row or column coset.
The coset is an affine $3$-space over $\F$, where further occupation
bounds allow at most three selected points in every affine plane.
These constraints, together with~\eqref{eq:intro-budget}, force the
matrix-rank profile
\begin{equation}\label{eq:intro-profile}
 (n_1,n_2,n_3)=(16,1,3),\qquad
 \sum_t\rank A_t=16+2+9=27.
\end{equation}

The last equality is decisive. The images of the $M_t$ span the
$27$-dimensional space, and their dimensions sum to $27$; they
therefore form a direct sum. Uniqueness of its components implies
\begin{equation}\label{eq:intro-saturation}
 M_tP^{-1}M_s=\delta_{ts}M_t.
\end{equation}
For the multiplication tensor, the product on the left can be computed
explicitly:
\begin{equation}\label{eq:intro-product}
 M(A,B,C)P^{-1}M(A',B',C')
   =M(AB'C^TA',B,C').
\end{equation}
The diagonal identities show that an invertible $A_t$ requires
invertible $B_t$ and $C_t$. An off-diagonal identity would then make a
product of invertible matrices zero if two first factors were
invertible. The required three in~\eqref{eq:intro-profile} are
impossible.

The resulting obstruction is useful separately from the finite
classification: in \emph{any} decomposition whose first-factor matrix
ranks sum to $27$, at most one first factor is invertible
(Proposition~\ref{prop:invertibility} and
Corollary~\ref{cor:profiles}). For a saturated $22$-term
decomposition with nonzero first factors, this leaves only the profiles
$(17,5,0)$ and $(18,3,1)$. Thus the equality argument gives both the
final contradiction at length $20$ and restrictions on possible shorter
algorithms beyond the known upper bound.

\subsection{Background and related methods}
Bl\"aser's substitution argument gives the classical lower bound
$19$ for $3\times3$ multiplication~\cite{blaser}.
On the constructive side, SAT methods have produced many inequivalent
$23$-term algorithms~\cite{heule}, while reinforcement learning has
extended the search for exact decompositions across matrix
formats~\cite{alphatensor}. These developments make small
matrix-multiplication tensors a meeting point of algebraic complexity,
computational search, and automated mathematical reasoning.

Wang's automated finite-field framework~\cite{wang20} raised the
binary lower bound to $20$ using symmetry classes of restrictions and
recursive lower-bound certificates. We use its earlier,
$496$-representative quotient catalogue with the numeric labels fixed.
The capacity viewpoint also appears in D'Ambrosio's determination of
the binary rank of $\langle2,3,4\rangle$~\cite{dambrosio}.

Wang's proof of $21$~\cite{wang21} and ours share the occupation
bounds and the affine rank-one coset reduction. His argument
strengthens a restriction table by rank-one-span searches and
backtracking, then excludes rank-one completions of the remaining
high-rank configurations by exhaustive enumeration. Here integer
occupation certificates establish eight additional quotient bounds;
affine-plane capacities force~\eqref{eq:intro-profile}, and
\eqref{eq:intro-saturation}--\eqref{eq:intro-product} give an algebraic
contradiction involving all three factor lists.
A public review draft in Tahir's repository~\cite{tahir} gives another
counting approach, using singular first factors and
branch-and-bound exclusions of seven normal forms.

The row/column-coset lemma is a special case of the matrix-adjacency
geometry initiated by Hua~\cite{hua}; we give an elementary proof.
The rank-additivity lemma is proved over an arbitrary field.
Their combination with the finite quotient bounds is what converts
first-factor occupation into the cross-factor obstruction for $\T$.

\subsection{Formal proof and research development}
The accompanying Lean development proves the tensor conventions,
the occupation and geometric arguments, the saturation identities, and
the finite lower bounds used by the proof. Its final theorem has no
assumed finite-bound premises. The development also proves the full
$496$-representative quotient catalogue and its transport to arbitrary
subspaces. The finite computations enter through exact integer branch
certificates whose leaves and exhaustive splits are checked within Lean.

Qiushi Engine developed the structural argument during sustained
autonomous research on exact matrix multiplication.
The mathematical progression was from shorter-algorithm searches to
quotient cores, from admissible first-factor supports to the question
of factor compatibility, and finally to equality in the full tensor's
rank bound. Section~\ref{sec:development} explains these transitions
and their role in the discovery. The accompanying research
trajectory~\cite{qiushi-report} connects them to the supporting notes,
programs, counterexamples, and results, preserving intermediate
knowledge as well as the final proof.

The exposition follows the proof. Section~\ref{sec:preliminaries} fixes
the tensor and quotient conventions; Section~\ref{sec:finite} states
and certifies the finite premises. Sections~\ref{sec:geometry}
and~\ref{sec:saturation} establish the profile and its contradiction.
The remaining sections give consequences, the formalization, and the
research development. The appendices specify the finite data, a
$23$-term witness, and the formal interfaces and verification record.
